\documentclass[12pt,a4paper]{article}

%% PACHETE MATEMATICE
\usepackage{amsmath,amssymb,amsfonts}
\usepackage{amsthm}
\usepackage{mathpazo}
\usepackage{eulervm}
\usepackage{enumerate}
\usepackage{color}
\usepackage{graphicx}
\usepackage[all]{xy}
\usepackage{xcolor}
\usepackage{framed}
\usepackage[unicode]{hyperref}
\setcounter{MaxMatrixCols}{20}
\usepackage{multicol}
% \usepackage[romanian]{babel}


%% ASPECT
\linespread{1.05}
\usepackage[T1]{fontenc}
\usepackage[utf8x]{inputenc}
\usepackage[margin=2cm]{geometry}
% \usepackage[color]{showkeys}
\usepackage{anyfontsize}
\usepackage{titlesec}
\titleformat*{\section}{\large\bfseries}

\usepackage{fancyhdr}
\pagestyle{fancy}
\rhead{\small \color{gray} Notițe de Adrian Manea}
\lhead{\small \color{gray} ETTI-AM2, 2017---2018, M. Joița \& A. Niță}


\usepackage[autostyle = false, style = german]{csquotes}
\usepackage[romanian]{babel}
\newcommand\qq{\enquote}

%% COMENZILE MELE
\renewcommand*{\proofname}{Dem.:}
\newcommand\bb{\mathbb}
\newcommand\dr{\mathrm}
\newcommand\kal{\mathcal}
\newcommand\fk{\mathfrak}
\newcommand\op{\oplus}
\newcommand\ot{\otimes}
\newcommand\ds{\displaystyle}
\newcommand\id{\indent}
\newcommand\nid{\noindent}
\newcommand\seq{\subseteq}
\newcommand\sneq{\subsetneq}
\newcommand\speq{\supseteq}
\newcommand\spneq{\supsetneq}
\newcommand\rar{\rightarrow}
\newcommand\Rar{\Rightarrow}
\newcommand\xrar{\xrightarrow}
\newcommand\lar{\leftarrow}
\newcommand\Lar{\Leftarrow}
\newcommand\xlar{\xleftarrow}
\newcommand\lrar{\leftrightarrow}
\newcommand\Lrar{\Leftrightarrow}
\newcommand\ideal{\trianglelefteq}
\newcommand\ai{\text{ a.\^{i}. }}
\newcommand\sk{(\textit{Schi\c{t}\u{a}}) }
\newcommand\rhup{\rightharpoonup}
\newcommand\rhdn{\rightharpoondown}
\newcommand\lhup{\leftharpoonup}
\newcommand\lhdn{\leftharpoondown}
\newcommand\vid{\emptyset}
\newcommand\wh{\widehat}
\newcommand\ol{\overline}
\newcommand\NN{\mathbb{N}}
\newcommand\RR{\mathbb{R}}
\newcommand\ZZ{\mathbb{Z}}
\newcommand\QQ{\mathbb{Q}}
\newcommand\CC{\mathbb{C}}

%% STILURILE MELE
\newtheoremstyle{dotless}{}{}{\itshape}{}{\bfseries}{:}{ }{}

\theoremstyle{dotless}
\newtheorem{theorem}{Teorem\u{a}}
\newtheorem{proposition}{Propozi\c{t}ie}
\newtheorem{lemma}{Lem\u{a}}
\newtheorem{corollary}{Corolar}
\newtheorem{exercise}{Exerci\c{t}iu}

\newtheoremstyle{dotlessdr}{}{}{}{}{\bfseries}{:}{ }{}
\theoremstyle{dotlessdr}
\newtheorem{example}{Exemplu}
\newtheorem{definition}{Defini\c{t}ie}
\newtheorem{remark}{Observa\c{t}ie}




\begin{document}
% \definecolor{labelkey}{RGB}{222,0,0}

\begin{center}
  {\large\textbf{Seminar 9 \\
      Integrale complexe}}
\end{center}

\vspace{2cm}

\section*{Teorema lui Cauchy}

\indent\indent Definiția și proprietățile esențiale ale integralelor din cazul real
rămîn valabile și în cazul complex, motiv pentru care le vom omite, începînd cu elementele
de noutate. De fapt, dat fiind că numerele complexe impun o structură bidimensională
prin reprezentarea în plan, majoritatea calculelor de integrale complexe vor fi
similare cu integralele curbilinii.

De exemplu, două calcule simple:

\[
  I_1 = \ds\int_{|z| = 1} z |dz|.
\]

Scriind polar $ z = e^{it}, t \in [0, 2\pi] $, avem $ dz = ie^{it} dt $, deci
$ |dz| = dt $. Atunci:
\[
  I_1 = \int_0^{2\pi} e^{it} dt = \frac{1}{i} e^{it} \Big|_0^{2\pi} = 0.
\]

\[
  I_2 = \int_S z |dz|,
\]
unde $ S $ este segmentul care unește pe $ 0 $ și $ i $.

Putem scrie segmentul parametric $ z = ti, t \in [0, 1] $, deci $ dz = idt $ și
$ |dz| = dt $. Atunci:
\[
  I_2 = \int_0^1 tidt = \frac{i}{2}.
\]

O teoremă foarte utilă pentru calculul integralelor complexe în lungul unor
curbe speciale este următoarea.

\begin{theorem}[Cauchy]\label{thm:cauchy}
  Fie $ D \seq \CC $ un domeniu simplu conex și fie $ f : D \to \CC $ o funcție
  olomorfă pe $ D $, astfel încît $ P = \mathrm{Re} f, Q = \mathrm{Im} f $ să fie
  de clasă $ \kal{C}^1(D) $.

  Fie $ \gamma : [a, b] \to D $ o curbă închisă jordaniană\footnote{fără autointersecții}
  de clasă $ \kal{C}^1 $ pe porțiuni astfel încît compactul $ \mathrm{Int}\gamma $ să
  verifice condițiile formulei Green-Riemann, atunci $ \ds\int_\gamma f(z) dz = 0 $.
\end{theorem}

Teorema are loc, de fapt, pentru orice curbă $ \gamma $ închisă, jordaniană și
\emph{omotopă cu zero}, adică poate fi deformată continuu, fără rupere, la un punct.

În particular, avem:

\begin{corollary}\label{cor:cauchy}
  Fie $ f : D \to \CC $ o funcție olomorfă pe un domeniu simplu conex $ D \seq \CC $
  și fie $ \gamma_1, \gamma_2 $ două drumuri de clasă $ \kal{C}^1 $ avînd aceleași
  capete și situate în $ D $. Atunci:
  \[
    \int_{\gamma_1} f(z) dz = \int_{\gamma_2} f(z) dz.
  \]
\end{corollary}

De exemplu:
\[
  \int_{|z| = r} \frac{e^z \sin z}{1 - z^3} dz = 0,
\]
deoarece putem considera cercul $ C = \{z \in \CC \mid |z - a| = r > 0 \} $ ca pe
un drum închis jordanian, de clasă $ \kal{C}^1 $, orientat pozitiv, parcurs o
singură dată în sens trigonometric direct.

\vspace{1cm}

Următorul rezultat este esențial în calculele ce vor urma.

\begin{theorem}[Formula integrală Cauchy]\label{thm:cauchy-int}
  Fie $ D \seq \CC $ un domeniu și $ f : D \to \CC $ o funcție olomorfă pe $ D $.
  Fie $ \overline{\Delta} \seq D $, unde $ \Delta $ este un domeniu simplu conex,
  mărginit cu frontiera $ \gamma $, o curbă închisă jordaniană, de clasă $ \kal{C}^1 $
  pe porțiuni, orientată pozitiv.

  Atunci pentru orice punct $ a \in \Delta $ fixat are loc formula:
  \[
    f(a) = \frac{1}{2\pi i} \int_\gamma \frac{f(z)}{z - a} dz.
  \]
\end{theorem}

De exemplu, să calculăm $ \ds\int_{|z - 2i| = 1} \frac{1}{z^2 + 4} dz $.

Putem scrie integrala astfel:
\[
  \int_{|z - 2i| = 1} \frac{1}{z^2 + 4} dz = %
  \int_{|z - 2i| = 1} \frac{\frac{1}{z + 2i}}{z - 2i} dz = %
  \int_{|z - 2i| = 1} \frac{f(z)}{z - 2i} dz,
\]
unde am introdus $ f(z) = \frac{1}{z + 2i} $.

Aplicînd formula integrală a lui Cauchy, obținem:
\[
  f(2i) = \frac{1}{2 \pi i} \int_{|z - 2i| = 1} \frac{f(z)}{z - 2i} dz \Rightarrow %
  \int_{|z - 2i| = 1} \frac{f(z)}{z - 2i} dz = 2 \pi i f(2i) = \frac{\pi}{2}.
\]

%%%%%%%%%%%%%%%%%%%%%%%%%%%%%%%%%%%%%%%%%%%%%%%%%%%%%%%%%%%%%%%%%%%%%%

\section*{Funcții analitice. Serii Laurent}

Noțiunile care urmează sînt adaptate din cazul seriilor reale de puteri.

\begin{theorem}\label{thm:serie-p}
  Fie $ z_0 \in \CC $ fixat și $ S(z) = \ds\sum_{n \geq 0} a_n(z - z_0)^n, \forall z $,
  cu $ |z - z_0| < R $ suma seriei centrate în $ z_0 $, iar $ R $ este raza de
  convergență.

  Atunci funcția $ S $ este olomorfă în orice punct $ z \in B(z_0, R) $ și are loc:
  \[
    S'(z) = \sum_{n \geq 0} na_n(z - z_0)^{n - 1}, \forall z \in B(z_0, R).
  \]
\end{theorem}

În particular, rezultă:

\begin{corollary}\label{cor:serie-p}
  Fie $ z_0 \in \CC $ fixat și $ S(z) $ ca în teoremă.

  Atunci funcția $ S(z) $ are derivate complexe de orice ordin în orice punct
  $ z \in B(z_0, R) $ și $ a_k = \dfrac{S^{(k)}(z_0)}{k!}, \forall k \in \NN $.
\end{corollary}

Funcțiile complexe cu proprietăți \qq{bune} pentru analiză sînt cele \emph{analitice}.

\begin{definition}\label{def:f-analitica}
  Fie $ A \seq \CC $ o mulțime deschisă. Funcția $ f : A \to \CC $ se numește
  \emph{analitică} pe $ A $ dacă $ \forall z_0 \in A $, există o serie formală
  $\sum_{n \geq 0} a_n X^n $ convergentă într-un disc de rază $ R > 0 $,
  astfel încît există discul $ B(z_0, r) \seq A $, cu $ 0 < r \leq R $
  cu proprietatea:
  \[
    f(z) = \sum_{n \geq 0} a_n(z - z_0)^n, \forall z \in B(z_0, r).
  \]

  Cu alte cuvinte, funcțiile analitice pot fi dezvoltate în serie Laurent
  convergentă într-un disc de orice rază din domeniul de definiție.
\end{definition}

Utilitatea acestei noțiuni rezultă și din:

\begin{proposition}\label{pr:analitica-der}
  Fie $ A \seq \CC $ o mulțime deschisă și $ f : A \to \CC $ o funcție analitică
  pe $ A $. Atunci există derivatele complexe de orice ordin ale lui $ f $ în $ A $
  și într-o vecinătate a oricărui punct $ z_0 \in A $ avem:
  \[
    f(z) = \sum_{n \geq 0} \frac{f^{(n)}(z_0)}{n!} (z - z_0)^n.
  \]
\end{proposition}

Legătura dintre analiticitate și olomorfie este dată în următoarea teoremă.

\begin{theorem}[Weierstrass-Riemann-Cauchy]\label{thm:wrc}
  Fie $ A \seq \CC $ o mulțime deschisă și $ f : A \to \CC $ o funcție complexă.

  Atunci $ f $ este analitică pe $ A $ dacă și numai dacă $ f $ este olomorfă pe $ A $.
\end{theorem}

Seriile de puteri complexe se definesc mai jos.

\begin{definition}\label{def:laurent}
  Se numește \emph{serie Laurent centrată în $ z_0 \in \CC $} orice serie de funcții
  de forma
  \[
    \sum_{n \in \ZZ} a_n(z - z_0)^n, \quad a_n \in \CC.
  \]

  Seria se numește convergentă dacă \emph{partea principală} a ei, anume
  $ \sum_{n \geq 1} a_{-n}(z - z_0)^{-n} $ și \emph{partea Taylor}, anume
  $ \sum_{n \geq 0} a_n(z - z_0)^n $ sînt simultan convergente.
\end{definition}

%%%%%%%%%%%%%%%%%%%%%%%%%%%%%%%%%%%%%%%%%%%%%%%%%%%%%%%%%%%%%%%%%%%%%%

\section*{Puncte singulare. Reziduuri}

\indent\indent În continuare, vom studia punctele \qq{cu probleme} ale unei
funcții complexe și vom vedea cum putem face calcule analitice coerente care să le evite.

\begin{definition}\label{def:pc-izolat}
  Fie $ A \seq \CC $ o mulțime deschisă nevidă și $ f : A \to \CC $ o funcție
  olomorfă pe $ A $. Un punct $ z_0 \in \CC $ se numește \emph{punct singular izolat}
  al lui $ f $ dacă există un disc $ D = B(z_0, r > 0) $ astfel încît funcția să
  fie olomorfă pe discul punctat $ D - \{z_0\} $.
\end{definition}

De exemplu, punctul $ z = 2 $ este un punct singular izolat pentru funcția
$ f(z) = \frac{1}{z - 2} $.

\begin{definition}\label{def:pc-sing-ap}
  Fie $ f : A \to \CC $ o funcție olomorfă, cu $ A \seq \CC $ o mulțime deschisă
  nevidă și fie $ z_0 \in \CC $ un punct singular izolat al lui $ f $.

  Punctul $ z_0 $ se numește \emph{punct singular aparent} dacă seria Laurent
  $ f(z) = \sum_{n \in \ZZ } a_n(z - z_0)^n $ are partea principală nulă, adică
  $ a_n = 0, \forall n < 0 $.
\end{definition}

În fine:

\begin{definition}\label{def:pol}
  Punctul singular izolat $ z_0 $ se numește \emph{pol} dacă în seria Laurent asociată
  funcției $ f $, partea principală are un număr finit de termeni nenuli, număr care
  se numește \emph{ordinul polului}.

  Punctul singular izolat $ z_0 $ se numește \emph{punct singular esențial} dacă
  partea principală a seriei Laurent asociată funcției are o infinitate de termeni
  nenuli.
\end{definition}

În continuare, cîteva rezultate care ne ajută să identificăm natura punctelor
singulare pentru funcții complexe.

\begin{lemma}\label{le:sing-ap}
  Fie $ f : B(z_0; 0, r) \to \CC $ o funcție olomorfă pe coroana $ B(z_0; 0, r) $.

  Dacă $ |f(z)| \leq M $ pentru orice $ z \in B(z_0; 0, r) $, adică modulul funcției este
  mărginit, atunci $ z_0 $ este punct singular aparent al lui $ f $.
\end{lemma}

Mai mult, avem:

\begin{proposition}\label{pr:sing-ap}
  În condițiile și cu notațiile din lema de mai sus, $ z_0 $ este punct singular
  aparent pentru $ f $ dacă și numai dacă există și este finită limita:
  \[
    \ell = \lim_{z \to z_0} f(z).
  \]
\end{proposition}

De exemplu, $ z = 0 $ este singularitate aparentă pentru funcția $ f(z) = \frac{\sin z}{z} $,
deoarece, ca în cazul real, $ \lim_{z \to 0} f(z) = 1 $.

Mai departe, identificarea polilor se poate face cu următorul rezultat:

\begin{proposition}\label{pr:pol}
  Fie $ f : B(z_0; 0, r) \to \CC $ o funcție olomorfă pe coroana de definiție.

  Atunci punctul $ z_0 $ este pol pentru $ f $ dacă și numai dacă:
  \[
    \lim_{z \to z_0} f(z) = \infty.
  \]
\end{proposition}

De exemplu, considerînd funcția $ f(z) = \frac{z}{z - 1} $, observăm că $ z = 1 $
este pol simplu, deoarece $ \lim_{z \to 1} f(z) = \infty $, iar dezvoltarea în
serie Laurent a funcției $ f $ se scrie:
\[
  f(z) = \frac{1 + z - 1}{z - 1} = \frac{1}{z - 1} + 1,
\]
care are partea principală $ \frac{1}{z - 1} $. Așadar, avem un pol simplu (de ordinul 1).

Singularitățile esențiale se identifică astfel:

\begin{proposition}\label{pr:sing-es}
  Fie $ f : B(z_0; 0, r) \to \CC $ o funcție olomorfă pe coroana $ B(z_0; 0, r) $.

  Atunci punctul $ z_0 $ este singularitate esențială pentru $ f $ dacă și numai
  dacă nu există limita $ \lim_{z \to z_0} f(z) $.
\end{proposition}

De exemplu, pentru funcția $ f(z) = \exp\Big(\frac{1}{z}\Big) $, punctul $ z = 0 $
este singularitate esențială. Pe de o parte, putem vedea că limita din propoziție
nu există (depinde de direcția $ z \to 0 $), iar pe de alta, seria Laurent este:
\begin{align*}
  \exp\Big(\frac{1}{z}\Big) &= 1 + \frac{1}{1! z} + \frac{1}{2! z^2} + \dots \\
                            &= 1 + \frac{1}{1!} z^{-1} + \frac{1}{2!} z^{-2} + \dots,
\end{align*}
care are partea principală cu o infinitate de termeni nenuli.

\vspace{1cm}

Vom fi interesați de următorul obiect:

\begin{definition}\label{def:rez}
  Fie $ A \seq \CC $ o mulțime deschisă nevidă și fie $ f : A \to \CC $ o funcție
  olomorfă. Considerăm discul punctat (coroana) $ B(z_0; 0, r) \seq A $, astfel
  încît punctul $ z_0 $ să fie singularitate izolată a lui $ f $.

  Considerînd dezvoltarea în serie Laurent a funcției $ f $ în coroana $ B(z_0; 0,r) $
  sub forma:
  \[
    f(z) = \sum_{n \in \ZZ} a_n (z - z_0)^n,
  \]
  coeficientul $ a_{-1} $ se numește \emph{reziduul funcției} $ f $ în punctul
  singular $ z_0 $, notat $ \mathrm{Rez}(f, z_0) $.
\end{definition}

Reziduurile pot fi calculate în mai multe moduri, aparent indirecte.

\begin{proposition}\label{pr:rez}
  Păstrînd condițiile și notațiile din definiția de mai sus, fie un cerc
  de rază $ \rho $, anume $ C = \{ z \in \CC \mid |z - z_0| = \rho < r \} $, parcurs
  în sens trigonometric direct. Atunci:
  \[
    a_{-1} = \frac{1}{2 \pi i} \int_C f(z) dz.
  \]
\end{proposition}

De asemenea, avem și:

\begin{proposition}\label{pr:rez-mult}
  Fie $ z_0 \in \CC $ un pol de ordinul $ k > 0 $ pentru $ f $. Atunci:
  \[
    \mathrm{Rez}(f, z_0) = \frac{1}{(k - 1)!}\lim_{z \to z_0} %
    \big[ (z - z_0)^k f(z) \big]^{(k - 1)}.
  \]
\end{proposition}

În particular, obținem:

\begin{corollary}\label{cor:rez-mult}
  Fie $ f : B(z_0; 0, r) \to \CC $ o funcție olomorfă pe coroana de definiție, astfel
  încît se poate scrie $ f(z) = \frac{P(z)}{Q(z)}, \forall z \in B(z_0; 0, r) $, cu
  $ P, Q $ funcții olomorfe în discul $ B(z_0, r) $ și $ P(z_0) \neq 0, Q(z_0) = 0 $,
  iar $ Q'(z_0) \neq 0 $.

  Atunci $ z_0 \in \CC $ este un pol de ordinul întîi pentru $ f $ și:
  \[
    \mathrm{Rez}(f, z_0) = \frac{P(z_0)}{Q'(z_0)}.
  \]
\end{corollary}

Un caz particular este următorul:

\begin{definition}\label{def:rez-inf}
  Fie $ f : B(0; r, \infty) \to \CC $ o funcție olomorfă în exteriorul discului
  $ B(0, r) $, adică $ z = \infty $ este un punct singular izolat pentru $ f $.

  Se numește \emph{reziduul funcției $ f $ în punctul $ \infty $} reziduul funcției
  $ \big( -\frac{1}{t^2} \big) f\big(\frac{1}{t}\big) $ în punctul $ t = 0 $.
\end{definition}

Cel mai important rezultat al secțiunii este:

\begin{theorem}[Teorema reziduurilor]\label{thm:thm-rez}
  Fie $ D \seq \CC $ un domeniu și $ f : D - \{ \alpha_1, \dots, \alpha_k \} \to \CC $
  o funcție olomorfă pentru care $ \alpha_i $ sînt singularități izolate.

  Fie $ K \seq D $ un compact cu frontiera $ \Gamma = \partial K $, o curbă de clasă
  $ \kal{C}^1 $ pe porțiuni, jordaniană, orientată pozitiv și care conține toate $ \alpha_i $
  în interior. Atunci:

  \[
    \int_\Gamma f(z) dz = 2\pi i \sum_{j = 1}^k \mathrm{Rez}(f, \alpha_j).
  \]
\end{theorem}

În particular:

\begin{corollary}\label{cor:thrm-rez}
  În contextul și notațiile teoremei de mai sus, avem:
  \[
    \sum_{j = 1}^k \mathrm{Rez}(f, \alpha_j) + \mathrm{Rez}(f, \infty) = 0.
  \]
\end{corollary}

\vspace{1cm}

De exemplu, să calculăm integrala:
\[
  I = \int_{|z| = r} \frac{e^z}{(z - i)(z - 2)} dz, r > 0, r \neq 1, 2.
\]

\emph{Soluție:} Dacă $ 0 < r < 1 $, putem aplica teorema lui Cauchy (~\ref{thm:cauchy})
și găsim $ I = 0 $.

Dacă $ 1 < r < 2 $, aplicăm formula integrală a lui Cauchy (~\ref{thm:cauchy-int}) și găsim:
\[
  \int_{|z| = r} \frac{e^z}{(z - i)(z - 2)} dz = %
  \int_{|z| = r} \frac{\frac{e^z}{z - 2}}{z - i} dz = %
  \int_{|z| = r} \frac{f(z)}{z - i} dz,
\]
unde $ f(z) = \frac{e^z}{z - 2} $. Rezultă:
\[
  \int_{|z| = r} \frac{f(z)}{z - i} dz = 2 \pi i f(i) = 2 \pi i \frac{e^i}{i - 2}.
\]

Dacă $ r > 2 $, aplicăm teorema reziduurilor, cu $ i $ și $ 2 $ poli simpli. Avem:
\begin{align*}
  \mathrm{Rez}(f, i) &= \lim_{z \to i} (z - i) \frac{e^z}{(z - i)(z - 2)} = \frac{e^i}{i - 2} \\
  \mathrm{Rez}(f, 2) &= \lim_{z \to 2} \frac{e^z}{(z - i)(z - 2)} = \frac{e^2}{2 - i}.
\end{align*}

Rezultă, din teorema reziduurilor:
\[
  \int_{|z| = r} \frac{e^z}{(z - i)(z - 2)} dz = 2 \pi i \Big( \frac{e^i}{i - 2} + %
  \frac{e^2}{2 - i} \Big).
\]

%%%%%%%%%%%%%%%%%%%%%%%%%%%%%%%%%%%%%%%%%%%%%%%%%%%%%%%%%%%%%%%%%%%%%%

\section*{Exerciții}

1. Să se calculeze integrala $ \ds\int_\Gamma z^2 dz $, unde:
\begin{enumerate}[(a)]
\item $ \Gamma = [-1, i] \cup [i, 1] $;
\item $ \Gamma = \{ z(t) = 2 + it^2, 0 \leq t \leq 1 \} $;
\item $ \Gamma = \Big\{ z(t) = t + i \cos \frac{\pi t}{2}, -1 \leq t \leq 1 \Big\} $.
\end{enumerate}

\vspace{1cm}

2. Să se calculeze integrala $ \ds\int_\Gamma (z - a)^n dz $, unde $ \Gamma $ este
un cerc cu centrul în $ a $ și rază $ r > 0 $, iar $ n \in \ZZ $.

\vspace{1cm}

3. Folosind funcția $ f(z) = z, z \in \CC $ să se calculeze o integrală curbilinie
în două variante:
\begin{enumerate}[(a)]
\item $ I_1 = \ds\int_{AB} zdz $, unde drumul de integrare este sfertul de
  elipsă $ \frac{x^2}{a^2} + \frac{y^2}{b^2} = 1 $ din primul cadran,
  iar $ A(a, 0) $ și $ B(0, b) $;
\item $ I_2 = \ds\int_{AOB} zdz $, unde drumul de integrare se obține mergînd
  pe axele de coordonate.
\end{enumerate}

\vspace{1cm}

4. Să se calculeze integrala $ I = \ds\int_C z^2 dz $, unde $ C $ este segmentul
care unește originea $ O(0, 0) $ cu punctul $ A(2, 1) $.

\vspace{1cm}

5. Arătați că punctul $ z = a $ este singularitate aparentă pentru funcțiile:
\begin{enumerate}[(a)]
\item $ f : \CC - \{ z \mid z^3 + 1 = 0 \} \to \CC, f(z) = \frac{z^2 - 1}{z^3 + 1} $,
  cu $ a = -1 $;
\item $ g : \CC^\ast \to \CC, g(z) = \frac{\sin z}{z}, a = 0 $;
\item $ h : \CC^\ast \to \CC, h(z) = \frac{z}{\tan z}, a = 0 $;
\item $ k : \CC^\ast \to \CC, k(z) = \frac{\cos z - 1}{z}, a = 0 $.
\end{enumerate}

\vspace{1cm}

6. Arătați că punctul $ z = a $ este pol pentru funcțiile:
\begin{enumerate}[(a)]
\item $ f(z) = \frac{1}{z - 2}, a = 2 $;
\item $ f(z) = \frac{1}{(z^2 + 1)^2}, a = - i $;
\item $ f(z) = \frac{\cos z}{z}, a = 0 $;
\item $ f(z) = \frac{z^2 + 1}{z(z - 1)}, a = 0$;
\item $ f(z) = \frac{1}{1 - \cos z}, a = 0 $;
\item $ f(z) = \frac{\sin z}{z^3}, a = 0 $;
\item $ f(z) = \frac{1}{(z + 1)(z - 1)^3}, a = \pm 1 $.
\end{enumerate}

\vspace{1cm}

7. Calculați reziduurile funcțiilor în punctele $ a $ indicate:
\begin{enumerate}[(a)]
\item $ f(z) = \frac{\exp\big(z^2\big)}{z - 1}, a = 1 $;
\item $ f(z) = \frac{\exp\big(z^2\big)}{(z - 1)^2}, a = 1 $;
\item $ f(z) = \frac{z + 2}{z^2 - 2z}, a = 0 $;
\item $ f(z) = \frac{1 + e^z}{z^4}, a = 0 $;
\item $ f(z) = \frac{\sin z}{4z^2}, a = 0 $;
\item $ f(z) = \frac{z}{1 - \cos z}, a = 0 $.
\end{enumerate}

\end{document}